% Pandoc/Quarto template for den-paper extension
\documentclass[
  manuscript=article,
  layout=brief,
  year=2026,
  volume=1,
]{den}

% DOI

% Conference (for proceedings/poster)

% Dates
\published{20 May 2026}

% Editor and reviewers

% Strip abstract fields from bibliography entries (prevents LuaLaTeX parsing issues)

% Title
\title{User Guide for the DEN Paper Template}

% Authors and affiliations
\author{Meizahra Afidatie \orcid{true}}
% Running head uses the first author's family name only ("Gupta et al.").
% The \ifdefined guard makes only the first loop pass take effect.
\ifdefined\denshortauthorset\else
  \gdef\denshortauthorset{}%
  \shortauthor{Afidatie}%
\fi
\affiliation{Dewan Ekonomi Nasional}
\email{correspondence@email.domain}
\author{Muhamad Haikal \orcid{true}}
% Running head uses the first author's family name only ("Gupta et al.").
% The \ifdefined guard makes only the first loop pass take effect.
\ifdefined\denshortauthorset\else
  \gdef\denshortauthorset{}%
  \shortauthor{Haikal}%
\fi
\affiliation{Institution-2, City, Country}

% Keywords
\keywords{Ekspektasi; keyword-two; keyword number three}

% JEL classification codes
\jel{C22, O47}

% Abbreviations

% Pandoc/Quarto compatibility macros
\providecommand{\tightlist}{%
  \setlength{\itemsep}{0pt}\setlength{\parskip}{0pt}}
\providecommand{\citeproc}[2]{#1}
\providecommand{\citeproctext}{}

% Syntax highlighting macros for code blocks (defines the Shaded
% environment and the token commands Pandoc emits). Self-contained:
% it pulls in fancyvrb and the colour definitions it needs.

% Support for graphics
\RequirePackage{graphicx}

% Support for ding symbols (checkmarks, etc.)
\RequirePackage{pifont}

% Support for table column width calculations (required by Quarto-generated tables)
\RequirePackage{calc}

% Pandoc image handling - required for newer Pandoc/Quarto versions
% Constrains images to text width to prevent overflow
\makeatletter
\providecommand{\pandocbounded}[1]{%
  \sbox0{#1}%
  \ifdim\wd0>\linewidth
    \resizebox{\linewidth}{!}{#1}%
  \else
    #1%
  \fi
}
\makeatother

% CSL references support
\newlength{\cslhangindent}
\setlength{\cslhangindent}{1.5em}
\newlength{\csllabelwidth}
\setlength{\csllabelwidth}{3em}
\makeatletter
\newcommand{\CSLNoBibLabel}{\renewcommand{\@biblabel}[1]{}}
\makeatother
\newenvironment{CSLReferences}[2]
 {\begin{list}{}%
  {\setlength{\leftmargin}{0pt}
   \setlength{\itemindent}{0pt}
   \setlength{\itemsep}{#2\baselineskip}
   \setlength{\parsep}{0pt}
   \setlength{\topsep}{0pt}
   \setlength{\partopsep}{0pt}
   \ifodd #1
     \setlength{\leftmargin}{\cslhangindent}
     \setlength{\itemindent}{-\cslhangindent}
   \fi}
  \CSLNoBibLabel}
 {\end{list}}

% Header includes
\makeatletter
\@ifpackageloaded{caption}{}{\usepackage{caption}}
\AtBeginDocument{%
\ifdefined\contentsname
  \renewcommand*\contentsname{Daftar Isi}
\else
  \newcommand\contentsname{Daftar Isi}
\fi
\ifdefined\listfigurename
  \renewcommand*\listfigurename{Daftar Gambar}
\else
  \newcommand\listfigurename{Daftar Gambar}
\fi
\ifdefined\listtablename
  \renewcommand*\listtablename{Daftar Tabel}
\else
  \newcommand\listtablename{Daftar Tabel}
\fi
\ifdefined\figurename
  \renewcommand*\figurename{Gambar}
\else
  \newcommand\figurename{Gambar}
\fi
\ifdefined\tablename
  \renewcommand*\tablename{Tabel}
\else
  \newcommand\tablename{Tabel}
\fi
}
\@ifpackageloaded{float}{}{\usepackage{float}}
\floatstyle{ruled}
\@ifundefined{c@chapter}{\newfloat{codelisting}{h}{lop}}{\newfloat{codelisting}{h}{lop}[chapter]}
\floatname{codelisting}{Daftar}
\newcommand*\listoflistings{\listof{codelisting}{Daftar Daftar}}
\makeatother
\makeatletter
\makeatother
\makeatletter
\@ifpackageloaded{caption}{}{\usepackage{caption}}
\@ifpackageloaded{subcaption}{}{\usepackage{subcaption}}
\makeatother

% Keep internal links clickable without PDF border underlines.
\hypersetup{pdfborder={0 0 0}}

% Source note environment for figure attributions (right-aligned, compact)
% \FloatBarrier forces pending floats to appear before the note, so the
% figure and its source line always land on the same page.
\RequirePackage{placeins}
\newenvironment{source}{%
  \FloatBarrier%
  \vspace{-1em}%
  \begin{flushright}\scriptsize\itshape
}{%
  \end{flushright}%
}
\newcommand{\sumber}[1]{%
  \FloatBarrier%
  \vspace{-1em}%
  \begin{flushright}\scriptsize\itshape
  Sumber: #1
  \end{flushright}%
}

% Allow URLs in bibliography to break across lines at any character.
% \UrlOrds adds break opportunities after all ordinary characters so long
% URLs (DOIs, web links) do not overflow the right margin.
\makeatletter
\g@addto@macro{\UrlBreaks}{\UrlOrds}
\makeatother

% Allow LaTeX more flexibility in line breaking to prevent overflow
% (needed for Indonesian/non-English text on narrow B5 paper)
\tolerance=1000
\emergencystretch=1.5em

\begin{document}


\section{Ekspektasi dan Kredibilitas dalam Kebijakan
Moneter}\label{ekspektasi-dan-kredibilitas-dalam-kebijakan-moneter}

\emph{Hall, Woodford, dan pelajaran bagi Indonesia}

Robert E. Hall dan Michael Woodford menerima Sveriges Riksbank Prize in
Economic Sciences in Memory of Alfred Nobel 2026 atas kontribusi mereka
dalam membangun fondasi mikroekonomi bagi kebijakan stabilisasi
makroekonomi. Hall memperlihatkan bahwa konsumsi rumah tangga lebih
tepat dipahami sebagai keputusan yang berorientasi ke depan, bukan
sekadar reaksi atas pendapatan hari ini. Sementara itu, Woodford
mengembangkan kerangka teoretis yang menghubungkan keputusan antarwaktu
rumah tangga, kekakuan penyesuaian harga perusahaan, dan kebijakan suku
bunga bank sentral.

Kerangka ini menjelaskan bagaimana kebijakan moneter memengaruhi inflasi
dan aktivitas ekonomi melalui ekspektasi pelaku ekonomi. Benang merah
keduanya adalah perilaku yang berorientasi ke depan
(\emph{forward-looking}). Artinya, dinamika ekonomi agregat berakar pada
keputusan jutaan rumah tangga dan perusahaan yang tidak hanya merespons
kondisi saat ini, tetapi juga mempertimbangkan perkembangan ekonomi dan
kebijakan pada masa mendatang (Dizikes 2026). Dalam model New Keynesian,
logika ini tecermin pada \emph{Phillips curve} yang menempatkan
ekspektasi inflasi sebagai penentu utama inflasi hari ini (Woodford
2003).

Tulisan ini membahas kontribusi Hall dan Woodford, serta pelajaran yang
dapat diambil Indonesia. Bagian pertama membahas Hall dan implikasinya
bagi perilaku konsumsi rumah tangga. Bagian kedua membahas Woodford dan
model New Keynesian, yang menekankan pentingnya ekspektasi dan
kredibilitas bank sentral. Bagian terakhir membahas pelajaran bagi
Indonesia, termasuk pentingnya mengelola ekspektasi masyarakat dan
kredibilitas kebijakan fiskal dan moneter.

\section{Lahirnya model ekonomi New
Keynesian}\label{lahirnya-model-ekonomi-new-keynesian}

Hingga 1960-an, kebijakan makroekonomi dirancang dengan model Keynesian
tradisional yang dibangun dari hubungan statistik antarvariabel makro
dalam data. Salah satu hubungan terpentingnya adalah \emph{Phillips
curve}, yang menunjukkan bahwa inflasi tinggi biasanya disertai
pengangguran rendah. Berbekal hubungan ini, pembuat kebijakan merasa
bisa memilih kombinasi yang diinginkan, misalnya menerima inflasi
sedikit lebih tinggi agar pengangguran turun. Keyakinan itu terpatahkan
pada 1970-an, ketika inflasi dan pengangguran justru naik bersamaan.
Lucas kemudian menjelaskan bahwa hubungan semacam ini tidak dapat
diandalkan untuk merancang kebijakan karena masyarakat akan mengubah
ekspektasi dan perilakunya begitu kebijakan berubah (Lucas 1976). Kritik
ini mendorong lahirnya model \emph{Real Business Cycle}, yang dibangun
dari keputusan rasional rumah tangga dan perusahaan (Muth 1961). Namun,
model ini menganggap harga selalu menyesuaikan seketika sehingga
kebijakan moneter tidak berpengaruh terhadap output dan lapangan kerja,
padahal pengalaman bank sentral menunjukkan sebaliknya. Model New
Keynesian kemudian hadir sebagai jalan tengah. Model ini tetap berangkat
dari keputusan rumah tangga dan perusahaan yang memperhitungkan masa
depan, tetapi mengakui bahwa harga tidak berubah setiap saat sehingga
kebijakan moneter kembali punya peran.

Mengapa kabar mengenai kenaikan pendapatan di masa depan dapat
memengaruhi konsumsi seseorang hari ini, bahkan sebelum tambahan
pendapatan tersebut benar-benar diterima? Sebaliknya, mengapa kenaikan
pendapatan yang sudah diperkirakan sebelumnya belum tentu mengubah
konsumsi ketika akhirnya terealisasi? Pertanyaan ini berkaitan erat
dengan salah satu kontribusi terpenting Robert E. Hall dalam
makroekonomi, yaitu bagaimana ekspektasi dan informasi baru menentukan
perilaku konsumsi rumah tangga. Berangkat dari
\emph{Life-Cycle--Permanent Income Hypothesis} (LC--PIH) yang
dikembangkan melalui pemikiran Modigliani dan Friedman (Modigliani dan
Brumberg 1954; Friedman 1957), Hall (1978) mengubah cara ekonom memahami
sekaligus menguji keputusan konsumsi.

Dalam kerangka LC--PIH, konsumsi tidak hanya ditentukan oleh pendapatan
saat ini, tetapi juga oleh ekspektasi mengenai sumber daya yang tersedia
sepanjang hidup atau \emph{lifetime resources}. Oleh karena itu,
kenaikan pendapatan sementara tidak selalu diikuti peningkatan konsumsi
dalam jumlah yang sama. Sebaliknya, informasi mengenai kenaikan
pendapatan yang bersifat permanen dapat mendorong rumah tangga
meningkatkan konsumsi bahkan sebelum pendapatan tambahan tersebut
diterima, karena informasi tersebut mengubah ekspektasi mereka terhadap
pendapatan di masa depan.

Hall (1978) mengembangkan gagasan ini dengan menggabungkan LC--PIH dan
\emph{rational expectations}. Melalui \emph{Euler equation}, yang
menggambarkan kondisi optimal dalam keputusan konsumsi antarperiode,
Hall menunjukkan bahwa dengan sejumlah asumsi tertentu, termasuk
\emph{quadratic utility} dan hubungan antara tingkat bunga riil dan
\emph{rate of time preference}, konsumsi dapat mengikuti pola
\emph{random walk}. Implikasi utamanya adalah perubahan konsumsi pada
periode berikutnya seharusnya hanya mencerminkan informasi baru yang
sebelumnya tidak dapat diprediksi. Dengan demikian, informasi yang telah
tersedia pada periode sebelumnya tidak seharusnya memiliki
\emph{additional predictive power} terhadap perubahan konsumsi. Prediksi
inilah yang kemudian menjadi dasar pengujian empiris Hall terhadap
LC--PIH dengan \emph{rational expectations}.

\[
C_{t+1}=C_t+\varepsilon_{t+1}, \qquad E_t[\varepsilon_{t+1}]=0.
\]

Prediksi tersebut bukan berarti konsumsi bergerak secara acak tanpa
adanya pertimbangan ekonomi. Justru sebaliknya, apabila rumah tangga
telah menggunakan seluruh informasi yang tersedia untuk menentukan
konsumsi hari ini, perubahan konsumsi pada periode berikutnya seharusnya
hanya mencerminkan informasi baru yang sebelumnya tidak dapat
diprediksi. Sebagai ilustrasi, seseorang mengetahui pada Januari bahwa
gajinya akan meningkat secara permanen mulai Februari. Konsumsinya dapat
meningkat sejak informasi tersebut diterima dan dipercaya, tanpa harus
menunggu kenaikan gaji terealisasi. Ketika Februari tiba, kenaikan gaji
tersebut tidak lagi memberikan informasi baru sehingga tidak ada alasan
tambahan untuk menyesuaikan konsumsi. Namun, apabila kenaikan gaji
dibatalkan secara tidak terduga, konsumsi akan kembali disesuaikan untuk
merespons perubahan ekspektasi tersebut.

Terobosan Hall tidak hanya terletak pada prediksi \emph{random walk},
tetapi juga pada pendekatan empiris yang memungkinkan teori konsumsi
diuji tanpa harus mengetahui seluruh informasi yang dimiliki rumah
tangga. Hall (1978) menunjukkan bahwa apabila LC--PIH dengan
\emph{rational expectations} berlaku, informasi yang sudah tersedia pada
periode sebelumnya seharusnya tidak memiliki \emph{additional predictive
power} terhadap perubahan konsumsi. Implikasi ini dapat diuji melalui
regresi sederhana berikut.

\begin{equation}\protect\phantomsection\label{eq-hall1}{
\Delta C_{t+1}=a+bX_t+u_{t+1}, \qquad H_0=0.
}\end{equation}

Dalam persamaan Persamaan~\ref{eq-hall1}, (X\_t) mewakili informasi yang
sudah tersedia pada periode (t). Jika hipotesis nol berlaku (b=0),
informasi tersebut tidak memberikan kemampuan tambahan untuk memprediksi
perubahan konsumsi berikutnya. Pendekatan ini membedakan kontribusi Hall
dari model konsumsi sebelumnya yang lebih mengandalkan hubungan historis
antara pendapatan dan pengeluaran. Alih-alih menganggap hubungan
tersebut selalu berlaku, Hall menguji prediksi yang diturunkan dari
keputusan rumah tangga yang mempertimbangkan kondisi di masa depan.
Pendekatan ini juga sejalan dengan persoalan yang diangkat dalam
\emph{Lucas critique} (Lucas 1976), bahwa hubungan empiris yang diamati
pada masa lalu belum tentu tetap berlaku ketika perubahan kebijakan
turut mengubah ekspektasi dan perilaku pelaku ekonomi.

Dalam pengujian empirisnya, Hall (1978) menggunakan sejumlah variabel,
termasuk pendapatan disposabel, konsumsi pada periode sebelumnya, dan
harga saham, untuk menguji apakah informasi masa lalu masih dapat
memprediksi perubahan konsumsi. Hasilnya menunjukkan bahwa pendapatan
pada periode sebelumnya tidak memiliki \emph{additional predictive
power} yang signifikan setelah konsumsi masa lalu diperhitungkan,
sejalan dengan prediksi \emph{random walk}. Namun, Hall menemukan bahwa
harga saham pada periode sebelumnya masih memiliki sejumlah
\emph{predictive power} terhadap perubahan konsumsi berikutnya. Temuan
ini menarik karena menunjukkan adanya penyimpangan dari prediksi
\emph{random walk}. Padahal, dalam kerangka teori tersebut, informasi
yang tercermin dalam harga saham seharusnya sudah diperhitungkan oleh
rumah tangga ketika menentukan konsumsi, sehingga tidak lagi memiliki
kemampuan untuk memprediksi perubahan konsumsi di masa depan.

Hall (1978) sendiri menekankan bahwa hasil \emph{random walk} bergantung
pada asumsi yang cukup ketat. Dalam penelitian berikutnya,
\emph{Intertemporal Substitution in Consumption}, Hall (1988) membahas
bagaimana rumah tangga mengalokasikan konsumsi antarperiode ketika
ekspektasi terhadap tingkat bunga riil berubah. Pembahasan ini berkaitan
dengan \emph{elasticity of intertemporal substitution} (EIS), yang
mengukur sensitivitas pertumbuhan konsumsi terhadap perubahan insentif
untuk menabung. Secara teori, kenaikan \emph{expected real return} akan
mendorong rumah tangga menunda sebagian konsumsi saat ini, tetapi
besarnya respons tersebut bergantung pada preferensi rumah tangga.

Dengan menggunakan data konsumsi dan \emph{asset returns} di Amerika
Serikat, Hall (1988) menemukan bukti yang relatif lemah mengenai
\emph{intertemporal substitution}. Estimasinya menunjukkan bahwa EIS
cenderung rendah, bahkan berpotensi mendekati nol dan kemungkinan tidak
lebih besar dari 0,2. Temuan ini mempertanyakan seberapa besar perubahan
konsumsi yang dapat dihasilkan oleh perubahan tingkat bunga riil yang
relatif kecil. Pada saat yang sama, hasil tersebut menunjukkan
pentingnya pengujian empiris yang lebih hati-hati, terutama ketika
ekspektasi, pengukuran, dan waktu pengambilan keputusan membuat hubungan
yang diprediksi teori tidak selalu terlihat jelas dalam data.

Cochrane (2024) menawarkan dua kemungkinan penjelasan atas temuan Hall
bahwa harga saham masih memiliki \emph{predictive power} terhadap
konsumsi. Pertama, harga saham tidak hanya mencerminkan ekspektasi
pendapatan di masa depan, tetapi juga perubahan \emph{expected returns}.
Kedua, adanya persoalan \emph{time aggregation dalam menghitung konsumsi
yang biasanya merupakan rata-rata di suatu periode tertentu}. Oleh
karena itu, harga saham merespons informasi baru lebih cepat
dibandingkan perubahan yang tercermin dalam data konsumsi bulanan atau
triwulanan. Akibatnya, harga saham dapat terlihat seolah memprediksi
perubahan konsumsi pada periode berikutnya, meskipun rumah tangga
sebenarnya telah menyesuaikan konsumsinya ketika informasi tersebut
muncul. Kedua kemungkinan ini menunjukkan bahwa penyimpangan dari
prediksi \emph{random walk} tidak serta-merta berarti rumah tangga gagal
membentuk \emph{rational expectations}.

Hubungan antara pasar saham dan konsumsi juga relevan dalam konteks
Indonesia, baik melalui \emph{direct wealth effect}, ketika perubahan
harga saham memengaruhi kekayaan pemilik aset keuangan, maupun melalui
\emph{sentiment channel}, ketika pergerakan pasar memengaruhi ekspektasi
rumah tangga terhadap kondisi ekonomi. Mansur, Gupta, dan Ginting (2026)
meneliti hubungan tersebut menggunakan data Indonesia periode
2019--2025, dengan \emph{Retail Sales Index} sebagai estimasi konsumsi
bulanan. Penelitian ini menggabungkan estimasi \emph{stock market
shocks} menggunakan GARCH(1,1) dengan analisis \emph{cointegration},
\emph{error correction model} (ECM), \emph{vector error correction
model} (VECM), dan \emph{impulse response functions}. Selain menguji
hubungan antara pergerakan harga saham dan konsumsi, penelitian ini juga
menganalisis waktu munculnya respons konsumsi serta perbedaan respons
terhadap \emph{market crashes} dan \emph{market booms}.

Hasil penelitian menunjukkan hubungan yang bersifat nonlinier dan
asimetris. Guncangan negatif ekstrem di pasar saham, yang diilustrasikan
melalui skenario kejatuhan harga saham sekitar 8 persen, berkaitan
dengan penurunan pertumbuhan penjualan ritel sebesar 5,9 poin persentase
setelah jeda waktu dua bulan. Respons yang signifikan secara statistik
terutama ditemukan pada guncangan yang melebihi 1,5 standar deviasi.
Sementara itu, \emph{error correction coefficient} (ECM) sebesar −0,374
menunjukkan bahwa sekitar 37,4 persen penyimpangan dari keseimbangan
jangka panjang terkoreksi setiap bulan. Menariknya, hanya guncangan
negatif (\emph{market crashes) }di pasar saham yang menunjukkan hubungan
signifikan secara statistik, sementara guncangan positif (\emph{market
booms) }tidak.

Pola tersebut menunjukkan bahwa penurunan tajam harga saham berkaitan
dengan respons konsumsi yang lebih besar dibandingkan kenaikan harga
saham dengan besaran yang sebanding. Mansur, Gupta, dan Ginting (2026)
menginterpretasikan temuan ini sebagai pola yang konsisten dengan
\emph{loss aversion}, \emph{precautionary saving}, dan
\emph{sentiment-driven transmission}. Namun, ketiganya dapat
mencerminkan mekanisme yang berbeda. \emph{Direct wealth effect} terjadi
ketika rumah tangga mengurangi konsumsi akibat penurunan nilai
portofolio saham yang dimiliki. Sebaliknya, \emph{sentiment channel}
dapat memengaruhi rumah tangga yang tidak memiliki saham sekalipun,
apabila penurunan pasar dipandang sebagai sinyal memburuknya prospek
pendapatan atau pekerjaan di masa depan.

Perbedaan tersebut penting karena hubungan agregat antara \emph{stock
market shocks} dan penjualan ritel belum dapat menunjukkan mekanisme
mana yang paling dominan. Harga saham dan konsumsi mungkin sama-sama
merespons perubahan kondisi ekonomi yang lebih luas. Dengan demikian,
penurunan harga saham dapat terjadi bersamaan dengan melemahnya
ekspektasi pertumbuhan ekonomi yang secara terpisah juga menekan
konsumsi.

Selain itu, temuan di Indonesia juga memunculkan pertanyaan mengenai
waktu dan besarnya penyesuaian konsumsi. Apakah jeda dua bulan
menunjukkan bahwa rumah tangga membutuhkan waktu untuk merespons
informasi baru, atau sebagian keterlambatan tersebut muncul akibat
\emph{time aggregation}, keterbatasan pengukuran \emph{retail sales},
dan waktu pelaksanaan pembelian? Sebagaimana dijelaskan Cochrane (2024),
keterlambatan yang terlihat dalam data tidak selalu mencerminkan
keterlambatan rumah tangga dalam memperbarui ekspektasinya.

Selain itu, respons yang asimetris menunjukkan pola yang lebih kompleks
dibandingkan prediksi model Hall yang paling sederhana. Namun, pola
tersebut tidak dengan sendirinya membuktikan adanya \emph{loss
aversion}. Faktor seperti \emph{borrowing constraints},
\emph{precautionary motives}, perubahan persepsi terhadap \emph{income
risk}, serta perbedaan persistensi antara \emph{positive} dan
\emph{negative shocks} juga dapat menghasilkan respons yang tidak
simetris. Terlebih lagi, \emph{Retail Sales Index} hanya mencakup
sebagian konsumsi rumah tangga dan tidak sepenuhnya merepresentasikan
konsumsi jasa, sehingga hasilnya tidak dapat langsung disamakan dengan
pengujian konsumsi agregat dalam Hall (1978).

Secara keseluruhan, kontribusi Hall memberikan landasan penting untuk
memahami bagaimana ekspektasi dan informasi baru memengaruhi keputusan
konsumsi rumah tangga. Meskipun temuan empiris menunjukkan bahwa
perilaku konsumsi lebih kompleks dibandingkan prediksi model yang paling
sederhana, pemikiran Hall tetap relevan dalam memahami bagaimana
masyarakat merespons perubahan kondisi ekonomi maupun kebijakan. Salah
satu implikasi pentingnya adalah bahwa dampak kebijakan tidak hanya
bergantung pada implementasinya, tetapi juga pada bagaimana informasi
mengenai kebijakan tersebut membentuk ekspektasi masyarakat. Dengan
demikian, kontribusi Hall tidak hanya memperluas pemahaman teoretis
mengenai perilaku konsumsi, tetapi juga menegaskan pentingnya
mempertimbangkan ekspektasi pelaku ekonomi dalam perumusan dan
komunikasi kebijakan.

\section{Ekspektasi dan Kredibilitas dalam Kebijakan
Moneter}\label{ekspektasi-dan-kredibilitas-dalam-kebijakan-moneter-1}

\subsection{Bapak dari Model New
Keynesian}\label{bapak-dari-model-new-keynesian}

\subsubsection{Michael Woodford}\label{michael-woodford}

Woodford menyusun analisis kebijakan moneter yang dibangun dari
keputusan optimal rumah tangga, perusahaan, dan bank sentral, bukan dari
korelasi antarvariabel agregat (Woodford 2003). Menurutnya, keberhasilan
kebijakan moneter tidak terutama ditentukan oleh kendali atas suku bunga
\emph{overnight}, melainkan oleh kemampuan bank sentral membentuk
ekspektasi pasar tentang arah suku bunga, inflasi, dan pendapatan ke
depan (Woodford 2003, 15). Karena itu, kredibilitas bank sentral menjadi
syarat utama efektivitas kebijakan.

\subsubsection{Model tiga persamaan}\label{model-tiga-persamaan}

Model New Keynesian dasar terdiri atas tiga persamaan yang mewakili
rumah tangga, perusahaan, dan bank sentral. Model ini mengasumsikan
perekonomian tertutup, dan seluruh variabel dinyatakan sebagai selisih
dari kondisi normalnya, misalnya inflasi sebagai selisih terhadap
target.

Pertama, kurva IS dinamis menggambarkan keputusan belanja rumah tangga
yang membandingkan belanja sekarang dengan belanja pada periode
berikutnya.

\begin{equation}\protect\phantomsection\label{eq-is}{
x_t = E_t(x_{t+1}) - \sigma\left(i_t - E_t(\pi_{t+1}) - r_t^n\right).
}\end{equation}

Dalam Persamaan~\ref{eq-is}, \(x_t\) adalah \emph{output gap}, yaitu
selisih output aktual terhadap output natural; \(E_t\) ekspektasi
berdasarkan informasi saat ini; \(i_t\) suku bunga kebijakan;
\(\pi_{t+1}\) inflasi periode berikutnya; \(r_t^n\) bunga riil natural;
dan \(\sigma > 0\) kepekaan belanja terhadap bunga riil. Jika bunga
riil, yaitu \(i_t - E_t(\pi_{t+1})\), lebih tinggi daripada bunga
natural, rumah tangga menunda belanja dan \emph{output gap} turun.
Artinya, suku bunga yang tidak berubah bisa membuat kebijakan efektif
lebih longgar ketika ekspektasi inflasi naik.

Kedua, \emph{Phillips curve} New Keynesian menggambarkan pembentukan
inflasi. Karena perusahaan tidak dapat mengubah harga setiap saat, harga
hari ini ditetapkan dengan memperhitungkan kondisi mendatang.

\begin{equation}\protect\phantomsection\label{eq-nkpc}{
\pi_t = \beta E_t(\pi_{t+1}) + \kappa x_t + v_t.
}\end{equation}

Dalam Persamaan~\ref{eq-nkpc}, \(\pi_t\) adalah inflasi relatif terhadap
target; \(\beta\) faktor diskonto antara nol dan satu; \(\kappa > 0\)
kepekaan inflasi terhadap \emph{output gap}; dan \(v_t\) gangguan biaya
(\emph{cost push}) di luar tekanan permintaan. Inflasi hari ini dibentuk
oleh ekspektasi inflasi, tekanan permintaan, dan gangguan biaya.

Ketiga, aturan suku bunga menggambarkan respons sistematis bank sentral
terhadap inflasi dan \emph{output gap}.

\begin{equation}\protect\phantomsection\label{eq-taylor}{
i_t = \phi_\pi \pi_t + \phi_x x_t + \varepsilon_t^i.
}\end{equation}

Dalam Persamaan~\ref{eq-taylor}, \(\phi_\pi\) dan \(\phi_x\) adalah
koefisien respons terhadap inflasi dan \emph{output gap}, sedangkan
\(\varepsilon_t^i\) mencerminkan keputusan di luar pola sistematis.
Aturan ini merupakan acuan analitis, bukan rumus resmi. Syarat
pentingnya dikenal sebagai \emph{Taylor principle}, yaitu
\(\phi_\pi > 1\), sehingga suku bunga naik lebih besar daripada inflasi
dan bunga riil ikut naik (Taylor 1993; Woodford 2003, 252--254).

\subsubsection{Ekspektasi yang dapat memenuhi dirinya
sendiri}\label{ekspektasi-yang-dapat-memenuhi-dirinya-sendiri}

Jika bank sentral merespons inflasi terlalu lemah, ekspektasi inflasi
dapat terwujud dengan sendirinya (\emph{self-fulfilling inflation}).
Kenaikan ekspektasi inflasi menurunkan bunga riil sehingga mendorong
belanja, sementara perusahaan menaikkan harga lebih awal. Inflasi
akhirnya mengikuti keyakinan masyarakat, bukan kebijakan, dan Woodford
menyebut kondisi ini indeterminasi. \emph{Taylor principle} mencegahnya
karena bunga riil selalu naik ketika ekspektasi inflasi meningkat,
sehingga target inflasi menjadi satu-satunya jangkar yang dipercaya
(Woodford 2003, bb. 2 dan 4).

\subsubsection{Komunikasi sebagai alat
kebijakan}\label{komunikasi-sebagai-alat-kebijakan}

Rumah tangga dan investor tidak hanya melihat suku bunga hari ini,
tetapi juga perkiraan suku bunga ke depan selama masa pinjaman atau
tenor investasinya. Karena itu, komunikasi bank sentral yang kredibel
dapat mengubah keputusan mereka sebelum suku bunga berubah, sehingga
kenaikan suku bunga yang diperlukan menjadi lebih kecil (Woodford 2003,
16; 2005). Penelitian di Amerika Serikat juga menunjukkan harga aset
bereaksi terhadap sinyal arah kebijakan The Fed, bukan hanya terhadap
perubahan suku bunganya (Gürkaynak, Sack, dan Swanson 2004). Namun,
pengumuman bank sentral juga dapat dibaca sebagai sinyal kondisi ekonomi
(Jarociński dan Karadi 2020).

\section{Pelajaran bagi indon}\label{pelajaran-bagi-indon}

Salah satu implikasi penting dari pemikiran Hall (1978) adalah kebijakan
ekonomi dapat memengaruhi keputusan konsumsi masyarakat bahkan sebelum
kebijakan tersebut diterapkan. Dalam kerangka \emph{Permanent Income
Hypothesis} dengan \emph{rational expectations}, rumah tangga
mempertimbangkan informasi mengenai pendapatan dan pengeluaran di masa
depan ketika menentukan konsumsi saat ini. Sebagaimana ditunjukkan
melalui prediksi \emph{random walk} Hall, perubahan konsumsi seharusnya
mencerminkan informasi baru yang sebelumnya tidak dapat diprediksi,
sementara informasi yang telah diantisipasi seharusnya sudah
diperhitungkan dalam keputusan konsumsi sebelumnya. Implikasi ini dapat
diilustrasikan melalui rencana kenaikan PPN dari 11 menjadi 12 persen
pada Januari 2025, sebagaimana telah diatur dalam UU Nomor 7 Tahun 2021
tentang Harmonisasi Peraturan Perpajakan. Namun, menjelang
implementasinya, pemerintah menerbitkan PMK No.~131 Tahun 2024 yang
mengatur agar beban PPN efektif untuk barang dan jasa nonmewah tetap
setara 11 persen. Masyarakat yang sebelumnya mengantisipasi kenaikan
pajak mungkin telah menyesuaikan rencana konsumsinya sebelum kebijakan
berlaku. Ketika cakupan penerapan PPN kemudian diperjelas, informasi
baru tersebut berpotensi kembali mengubah ekspektasi dan keputusan
konsumsi.

Contoh lainnya adalah kebijakan subsidi BBM yang berkaitan dengan
ekspektasi masyarakat terhadap biaya hidup di masa depan. Di tengah
ketidakpastian harga energi global, pada April 2026, pemerintah
menegaskan komitmen untuk mempertahankan harga BBM bersubsidi hingga
akhir tahun (Kementerian Koordinator Bidang Perekonomian Republik
Indonesia 2026). Dalam kerangka Hall, pengaruh pengumuman tersebut
terhadap konsumsi bergantung pada sejauh mana informasi yang disampaikan
mengubah ekspektasi rumah tangga mengenai sumber daya yang tersedia di
masa depan. Apabila masyarakat sebelumnya memperkirakan harga BBM akan
meningkat, komitmen pemerintah yang dianggap kredibel dapat mengurangi
kekhawatiran terhadap kenaikan biaya hidup dan berpotensi memengaruhi
keputusan konsumsi saat ini. Sebaliknya, apabila stabilitas harga telah
diantisipasi, pengumuman tersebut tidak seharusnya menimbulkan
penyesuaian konsumsi tambahan. Namun, respons aktual juga bergantung
pada kondisi rumah tangga, termasuk keterbatasan likuiditas, serta
kepercayaan masyarakat terhadap kemampuan pemerintah mempertahankan
komitmennya di tengah perubahan harga minyak dunia dan kondisi fiskal.

Dari kedua kasus tersebut, terdapat catatan penting bagi perumusan
kebijakan ekonomi. Dalam kasus PPN, kejelasan mengenai cakupan dan waktu
penerapan kebijakan menjadi penting agar masyarakat dapat merencanakan
pengeluaran berdasarkan informasi yang tepat. Sementara itu, kebijakan
subsidi BBM memperlihatkan pentingnya kredibilitas komitmen pemerintah
dalam membentuk ekspektasi mengenai biaya hidup di masa depan. Oleh
karena itu, pemerintah perlu memastikan bahwa kebijakan dikomunikasikan
secara jelas dan cukup dini, disertai penjelasan mengenai ketentuan,
jangka waktu, serta kondisi yang dapat menyebabkan penyesuaian
kebijakan. Perubahan kebijakan tidak serta-merta mengurangi kredibilitas
pemerintah, tetapi ketidakjelasan informasi atau perubahan tanpa
penjelasan yang memadai dapat menimbulkan ketidakpastian bagi
masyarakat. Dengan demikian, salah satu pelajaran penting dari pemikiran
Hall adalah bahwa efektivitas kebijakan tidak hanya berkaitan dengan
substansi dan implementasinya, tetapi juga dengan bagaimana masyarakat
menerima, memahami, dan memperhitungkan informasi tersebut dalam
keputusan ekonominya. Mempertimbangkan pembentukan ekspektasi masyarakat
perlu menjadi bagian dari perancangan dan evaluasi kebijakan ekonomi,
bukan sekadar persoalan komunikasi publik.

\section*{Referensi}\label{referensi}
\addcontentsline{toc}{section}{Referensi}

\protect\phantomsection\label{refs}
\begin{CSLReferences}{1}{1}
\bibitem[\citeproctext]{ref-cochrane2024}
Cochrane, John H. 2024. \emph{Bob Hall and Consumption}. Makalah untuk
konferensi penghormatan Robert E. Hall, Hoover Institution, 22 November
2024. \url{https://www.grumpy-economist.com/p/bob-hall-and-consumption}.

\bibitem[\citeproctext]{ref-dizikes2026}
Dizikes, Peter. 2026. \emph{Economists Robert Hall PhD '67 and Michael
Woodford PhD '83 Share Nobel Prize}. MIT News.
\url{https://news.mit.edu/2026/robert-hall-michael-woodford-share-nobel-prize-economics-1012}.

\bibitem[\citeproctext]{ref-friedman1957}
Friedman, Milton. 1957. \emph{A Theory of the Consumption Function}.
Princeton: Princeton University Press.
\url{https://www.nber.org/books-and-chapters/theory-consumption-function}.

\bibitem[\citeproctext]{ref-gurkaynak2004}
Gürkaynak, Refet S., Brian Sack, dan Eric T. Swanson. 2004. \emph{Do
Actions Speak Louder Than Words? The Response of Asset Prices to
Monetary Policy Actions and Statements}. Finance and Economics
Discussion Series Nos. 2004-66. Board of Governors of the Federal
Reserve System.
\url{https://www.federalreserve.gov/econres/feds/do-actions-speak-louder-than-words-the-response-of-asset-prices-to-monetary-policy-actions-and-statements.htm}.

\bibitem[\citeproctext]{ref-hall1978}
Hall, Robert E. 1978. {``Stochastic Implications of the Life
Cycle--Permanent Income Hypothesis: Theory and Evidence.''}
\emph{Journal of Political Economy} 86 (6): 971--987.
\url{https://doi.org/10.1086/260724}.

\bibitem[\citeproctext]{ref-hall1988}
Hall, Robert E. 1988. {``Intertemporal Substitution in Consumption.''}
\emph{Journal of Political Economy} 96 (2): 339--357.
\url{https://doi.org/10.1086/261539}.

\bibitem[\citeproctext]{ref-jarocinski2020}
Jarociński, Marek, dan Peter Karadi. 2020. {``Deconstructing Monetary
Policy Surprises---The Role of Information Shocks.''} \emph{American
Economic Journal: Macroeconomics} 12 (2): 1--43.
\url{https://doi.org/10.1257/mac.20180090}.

\bibitem[\citeproctext]{ref-pmk131-2024}
Kementerian Keuangan Republik Indonesia. 2024. \emph{Peraturan Menteri
Keuangan Nomor 131 Tahun 2024 tentang Perlakuan Pajak Pertambahan Nilai
atas Impor Barang Kena Pajak, Penyerahan Barang Kena Pajak, Penyerahan
Jasa Kena Pajak, Pemanfaatan Barang Kena Pajak Tidak Berwujud dari Luar
Daerah Pabean di Dalam Daerah Pabean, dan Pemanfaatan Jasa Kena Pajak
dari Luar Daerah Pabean di Dalam Daerah Pabean}.
\url{https://www.pajak.go.id/index.php/id/peraturan/perlakuan-pajak-pertambahan-nilai-atas-impor-barang-kena-pajak-penyerahan-barang-kena}.

\bibitem[\citeproctext]{ref-kemenko2026}
Kementerian Koordinator Bidang Perekonomian Republik Indonesia. 2026.
\emph{Pemerintah Berikan Kepastian Ketersediaan Energi dengan Harga
Terjangkau bagi Seluruh Lapisan Masyarakat}. Siaran pers.
\url{https://www.ekon.go.id/publikasi/detail/6873/pemerintah-berikan-kepastian-ketersediaan-energi-dengan-harga-terjangkau-bagi-seluruh-lapisan-masyarakat}.

\bibitem[\citeproctext]{ref-lucas1976}
Lucas, Jr., Robert E. 1976. {``Econometric Policy Evaluation: A
Critique.''} \emph{Carnegie-Rochester Conference Series on Public
Policy} 1: 19--46. \url{https://doi.org/10.1016/S0167-2231(76)80003-6}.

\bibitem[\citeproctext]{ref-mansur2026}
Mansur, Alfan, Krisna Gupta, dan Timothy Ginting. 2026. \emph{The
Asymmetric Wealth Effect: Evidence from Stock Market Crashes and Retail
Consumption in Indonesia}. DEN Working Paper No. 5. Dewan Ekonomi
Nasional.

\bibitem[\citeproctext]{ref-modigliani1954}
Modigliani, Franco, dan Richard Brumberg. 1954. {``Utility Analysis and
the Consumption Function: An Interpretation of Cross-Section Data.''}
Dalam \emph{Post-Keynesian Economics}, disunting oleh Kenneth K.
Kurihara, 388--436. New Brunswick, NJ: Rutgers University Press.

\bibitem[\citeproctext]{ref-muth1961}
Muth, John F. 1961. {``Rational Expectations and the Theory of Price
Movements.''} \emph{Econometrica} 29 (3): 315--335.
\url{https://doi.org/10.2307/1909635}.

\bibitem[\citeproctext]{ref-uu7-2021}
Republik Indonesia. 2021. \emph{Undang-Undang Nomor 7 Tahun 2021 tentang
Harmonisasi Peraturan Perpajakan}.
\url{https://www.pajak.go.id/id/peraturan/harmonisasi-peraturan-perpajakan}.

\bibitem[\citeproctext]{ref-taylor1993}
Taylor, John B. 1993. {``Discretion versus Policy Rules in Practice.''}
\emph{Carnegie-Rochester Conference Series on Public Policy} 39:
195--214. \url{https://doi.org/10.1016/0167-2231(93)90009-L}.

\bibitem[\citeproctext]{ref-woodford2003}
Woodford, Michael. 2003. \emph{Interest and Prices: Foundations of a
Theory of Monetary Policy}. Princeton: Princeton University Press.

\bibitem[\citeproctext]{ref-woodford2005}
Woodford, Michael. 2005. \emph{Central Bank Communication and Policy
Effectiveness}. NBER Working Paper No. 11898. National Bureau of
Economic Research. \url{https://doi.org/10.3386/w11898}.

\end{CSLReferences}

\end{document}
